\expandafter\ifx\csname CAPETINCLUDEMODE\endcsname\relax
\documentclass[11pt,a4paper]{article}
% preamble.tex — packages et commandes communes a toutes les fiches CAPET SII IE
\input{glyphtounicode}
\pdfgentounicode=1
\usepackage{cmap}
\usepackage[T1]{fontenc}
\usepackage[utf8]{inputenc}
\usepackage[french]{babel}
\usepackage[a4paper, top=1.6cm, bottom=1.6cm, left=1.4cm, right=1.4cm]{geometry}
\setlength{\headheight}{14.5pt}
\usepackage{amsmath, amssymb, amsfonts}
\usepackage{siunitx}
\sisetup{locale=FR, output-decimal-marker={,}, inter-unit-product=\ensuremath{{}\cdot{}}}
\usepackage{array, booktabs, multirow, tabularx}
\usepackage{xcolor}
\usepackage[most, breakable]{tcolorbox}
\usepackage{multicol}
\usepackage{enumitem}
\usepackage{mathtools}
\usepackage{lmodern}
\usepackage{microtype}
\usepackage{fancyhdr}
\usepackage{tikz}
\usetikzlibrary{positioning,arrows.meta}
\usepackage{lastpage}

% ── Couleurs par domaine ──────────────────────────────────────────────
\definecolor{colElec}{RGB}{26, 82, 148}       % Électrotechnique — bleu foncé
\definecolor{colElecLight}{RGB}{210, 225, 245} % fond bleu clair
\definecolor{colAuto}{RGB}{34, 120, 66}        % Automatique — vert
\definecolor{colMeca}{RGB}{185, 90, 20}        % Mécanique — orange
\definecolor{colTherm}{RGB}{160, 40, 40}       % Thermique — rouge
\definecolor{colQT}{RGB}{232, 244, 233}        % fond question type (vert très léger)
\definecolor{colJury}{RGB}{255, 243, 205}      % fond jury (jaune ambre)
\definecolor{colJuryFrame}{RGB}{200, 130, 0}   % bordure jury
\definecolor{colGrayLight}{RGB}{245, 245, 247} % fond gris neutre
\definecolor{colSeparator}{RGB}{180, 180, 180} % gris séparateur

% ── Style tcolorbox de section principale ────────────────────────────
\tcbset{
  ficheSection/.style={
    breakable,
    enhanced,
    colback=#1!7!white,
    colframe=#1,
    colbacktitle=#1,
    coltitle=white,
    fonttitle=\bfseries\normalsize,
    left=4pt, right=4pt, top=3pt, bottom=4pt,
    before skip=6pt, after skip=6pt,
  }
}

% ── Style encadré question type ──────────────────────────────────────
\tcbset{
  questiontype/.style={
    breakable,
    enhanced,
    colback=colQT,
    colframe=colElec!55!black,
    colbacktitle=colElec!55!black,
    coltitle=white,
    fonttitle=\bfseries\small,
    left=4pt, right=4pt, top=3pt, bottom=3pt,
    before skip=5pt, after skip=5pt,
  }
}

% ── Style encadré jury ───────────────────────────────────────────────
\tcbset{
  jury/.style={
    breakable,
    enhanced,
    colback=colJury,
    colframe=colJuryFrame,
    fonttitle=\bfseries\normalsize,
    left=4pt, right=4pt, top=3pt, bottom=3pt,
  }
}

% ── Mise en page ─────────────────────────────────────────────────────
\pagestyle{fancy}
\fancyhf{}
\renewcommand{\headrulewidth}{0.4pt}
\renewcommand{\footrulewidth}{0.3pt}

\setlist[itemize]{leftmargin=1.2em, itemsep=1pt, topsep=2pt, parsep=0pt}
\setlist[enumerate]{leftmargin=1.2em, itemsep=3pt, topsep=2pt, parsep=0pt}

\setlength{\parskip}{2pt}
\setlength{\columnsep}{10pt}

\def\CAPETBEGINFILE{\begin{document}}
\def\CAPETENDFILE{\end{document}}
\else
\def\CAPETBEGINFILE{}
\def\CAPETENDFILE{}
\fi

\fancyhead[L]{\textbf{\textcolor{colElec}{T08}} - Redresseurs (monophase et triphase)}
\fancyhead[R]{\small\textit{CAPET SII IE - Session 2026}}
\fancyfoot[C]{\small\thepage\,/\,\pageref{LastPage}}

\CAPETBEGINFILE

% -----------------------------------------------------------------------
%  BLOC EN-TeTE
% -----------------------------------------------------------------------

\begin{tcolorbox}[
  enhanced, breakable,
  colback=colElec,
  colframe=colElec,
  left=6pt, right=6pt, top=5pt, bottom=5pt,
  before skip=0pt, after skip=4pt
]
  {\color{white}\bfseries\large T08 \quad Redresseurs (monophase et triphase)}\\[2pt]
  {\color{colElec!30!white}\small
    \textbf{Domaine :} electrotechnique - electronique de puissance \quad
    \textbf{Presence :} \textbf{11/13} sujets (2013-2026) - theme \emph{tres frequent}
  }\\[1pt]
  {\color{colElec!20!white}\footnotesize
    Sous-themes : PD2 (Graetz monophase) - P3 (triphase simple alt.) - PD6 (triphase double alt.) - facteur de puissance - harmoniques - redresseur commande (thyristors)
  }
\end{tcolorbox}

% -----------------------------------------------------------------------
%  1. RAPPELS DE COURS
% -----------------------------------------------------------------------

\begin{tcolorbox}[
  ficheSection=colElec,
  title={1. Rappels de cours},
]

% -- 1.1 Tensions moyennes ---------------------------------------------
\textbf{1.1 Tensions moyennes de sortie}

\smallskip
\begin{center}
\renewcommand{\arraystretch}{1.5}
\begin{tabular}{lllll}
\toprule
\textbf{Pont} & \textbf{Diodes} & \textbf{Formule} & \textbf{Valeur num.} & \textbf{Pulsations/T} \\
\midrule
PD2 (monophase) & 4 & $\bar{U} = \dfrac{2\sqrt{2}}{\pi} V_{eff}$ & $\approx 0{,}9\,V_{eff}$ & 2 \\[8pt]
P3 (triphase) & 3 & $\bar{U} = \dfrac{3\sqrt{6}}{2\pi} V_{eff,simple}$ & $\approx 1{,}17\,V_{s,simple}$ & 3 \\[8pt]
PD6 (triphase) & 6 & $\bar{U} = \dfrac{3\sqrt{2}}{\pi} V_{eff,comp}$ & $\approx 2{,}34\,V_{s,simple}$ & 6 \\[8pt]
PD6 compact & 6 & $\bar{U} = 0{,}95 \cdot U_M$ & $U_M$ = ampli. composees & 6 \\
\bottomrule
\end{tabular}
\end{center}

\medskip
\hrule\vspace{4pt}

% -- 1.2 Facteur de puissance -----------------------------------------
\textbf{1.2 Facteur de puissance}

\smallskip
Pour un redresseur - diodes - $q$ pulsations (courant non sinusoedal - $f_P \neq \cos\varphi$) :
\[
  f_P = \frac{\sqrt{q}}{\pi} \sin\!\left(\frac{\pi}{q}\right)
\]

\begin{center}
\renewcommand{\arraystretch}{1.35}
\begin{tabular}{llll}
\toprule
\textbf{Pont} & $q$ & $f_P$ & \textbf{Conformite CDC $f_P \geq 0{,}9$} \\
\midrule
PD2 & 2 & $\dfrac{\sqrt{2}}{\pi} \approx 0{,}45$ & Non \\[6pt]
P3  & 3 & $\dfrac{\sqrt{3}}{\pi}\sin\!\left(\dfrac{\pi}{3}\right) \approx 0{,}827$ & Non \\[6pt]
PD6 & 6 & $\dfrac{\sqrt{6}}{\pi}\sin\!\left(\dfrac{\pi}{6}\right) \approx 0{,}955$ & \textbf{Oui} \\
\bottomrule
\end{tabular}
\end{center}

\medskip
\hrule\vspace{4pt}

% -- 1.3 Harmoniques injectes -----------------------------------------
\textbf{1.3 Harmoniques cete reseau}

\smallskip
Les redresseurs - diodes injectent des harmoniques de courant cete reseau.
Rangs presents : $6k \pm 1$ pour PD6 ($k = 1, 2, \ldots$) - rangs 5, 7, 11, 13, \ldots
Pour PD2 : rangs $2k \pm 1$ - rangs 3, 5, 7, \ldots\\
\textit{Taux d'ondulation tension : PD6 = 4,2\,\% - P3 = 18,3\,\% - PD2 = 48,2\,\%}

\end{tcolorbox}

% -----------------------------------------------------------------------
%  2. FORMULAIRE SYNTHeTIQUE
% -----------------------------------------------------------------------

\begin{tcolorbox}[
  ficheSection=colElec,
  title={2. Formulaire synthetique},
]
\small
\begin{multicols}{2}
\renewcommand{\arraystretch}{1.45}
\begin{tabular}{@{}p{4.6cm}ll@{}}
\toprule
\textbf{Formule} & \textbf{Usage} & \textbf{U.} \\
\midrule
\multicolumn{3}{@{}l}{\footnotesize\textit{Tensions moyennes}} \\
$\bar{U}_{PD2} = 0{,}9\,V_{eff}$ & Monoph. & V \\
$\bar{U}_{P3} = 1{,}17\,V_{s,simple}$ & 3 diodes & V \\
$\bar{U}_{PD6} = 2{,}34\,V_{s,simple}$ & 6 diodes & V \\
$\bar{U}_{PD6} = 0{,}95\,U_M$ & Ampli comp. & V \\
$U_{s,comp} = U_M/\sqrt{2}$ & Valeur eff. & V \\
\midrule
\multicolumn{3}{@{}l}{\footnotesize\textit{Rapport de transfo (amont)}} \\
$m = U_{s,comp}/U_{1,comp}$ & Yy & - \\
$m = V_{s,simple}/V_{1,simple}$ & Yy equiv. & - \\
\bottomrule
\end{tabular}

\columnbreak

\renewcommand{\arraystretch}{1.45}
\begin{tabular}{@{}p{4.6cm}ll@{}}
\toprule
\textbf{Formule} & \textbf{Usage} & \textbf{U.} \\
\midrule
\multicolumn{3}{@{}l}{\footnotesize\textit{Facteur de puissance}} \\
$f_P = \frac{\sqrt{q}}{\pi}\sin\!\left(\frac{\pi}{q}\right)$ & General & - \\
$f_{P,PD2} \approx 0{,}45$ & PD2 & - \\
$f_{P,P3} \approx 0{,}827$ & P3 & - \\
$f_{P,PD6} \approx 0{,}955$ & PD6 & - \\
\midrule
\multicolumn{3}{@{}l}{\footnotesize\textit{Thyristor commande}} \\
$\bar{U}_\alpha = \bar{U}_0\cos\alpha$ & Pont commande & V \\
\midrule
\multicolumn{3}{@{}l}{\footnotesize\textit{Courant efficace diode P3}} \\
$I_{D,eff} = I_S/\sqrt{3}$ & Chaque diode & A \\
\bottomrule
\end{tabular}
\end{multicols}
\normalsize
\end{tcolorbox}

% -----------------------------------------------------------------------
%  3. QUESTIONS TYPES
% -----------------------------------------------------------------------

\begin{tcolorbox}[
  ficheSection=colElec,
  title={3. Questions types - extraites des sujets reels},
]

% -- QT1 -------------------------------------------------------------

% Mini-procedure de resolution type copie concours
\begin{tcolorbox}[questiontype,title={Procedure type copie concours (5 etapes)}]
\begin{enumerate}[leftmargin=1.5em]
  \item \textbf{Qualifier le modele} : hypotheses, domaine de validite, grandeurs d'entree/sortie.
  \item \textbf{Ecrire les lois} : equations de base avec symboles et unites SI.
  \item \textbf{Isoler l'inconnue} : expression litterale, puis application numerique tracee.
  \item \textbf{Controler l'ordre de grandeur} : coherence physique et unite finale.
  \item \textbf{Conclure concours} : phrase courte orientee decision technique.
\end{enumerate}
\end{tcolorbox}

% Visuel de methode : chemin de resolution CAPET
\begin{center}
\begin{tikzpicture}[font=\small, node distance=7mm and 6mm, >=Stealth]
  \node[draw, rounded corners, fill=colElec!8, inner sep=4pt] (n1) {Notions};
  \node[draw, rounded corners, fill=colElec!8, right=of n1, inner sep=4pt] (n2) {Methode};
  \node[draw, rounded corners, fill=colElec!8, right=of n2, inner sep=4pt] (n3) {Application annale};
  \node[draw, rounded corners, fill=colElec!8, right=of n3, inner sep=4pt] (n4) {Verification jury};
  \draw[->, thick, colElec] (n1) -- (n2);
  \draw[->, thick, colElec] (n2) -- (n3);
  \draw[->, thick, colElec] (n3) -- (n4);
\end{tikzpicture}
\end{center}
\begin{tcolorbox}[questiontype,
  title={QT1 - Facteur de puissance P3 vs PD6 - choix de structure \hfill\normalfont\footnotesize E1 sujet 2026}
]
\textbf{enonce.}
Un reseau triphase $230/400\,\mathrm{V}$ alimente un chargeur de batterie DC. Le cahier des charges impose une tension DC de $600\,\mathrm{V}$ et un facteur de puissance $f_P \geq 0{,}9$. Calculer $f_P$ pour un pont P3 ($q = 3$) et pour un pont PD6 ($q = 6$). Conclure sur la structure - retenir.

\tcbline

\textbf{Corrige.}

\textbf{Pont P3 ($q = 3$) :}
\[
  f_{P,P3} = \frac{\sqrt{3}}{\pi} \cdot \sin\!\left(\frac{\pi}{3}\right)
  = \frac{1{,}732}{3{,}1416} \times 0{,}866 \approx 0{,}827
\]
$f_{P,P3} = 0{,}827 < 0{,}9$ - \textbf{non conforme}.

\textbf{Pont PD6 ($q = 6$) :}
\[
  f_{P,PD6} = \frac{\sqrt{6}}{\pi} \cdot \sin\!\left(\frac{\pi}{6}\right)
  = \frac{2{,}449}{3{,}1416} \times 0{,}5 \approx 0{,}955
\]
$f_{P,PD6} = 0{,}955 > 0{,}9$ - \textbf{conforme}. La solution PD6 est retenue.
\end{tcolorbox}

\vspace{4pt}

% -- QT2 -------------------------------------------------------------

\begin{tcolorbox}[questiontype,
  title={QT2 - PD6 - 600\,V DC - calcul du rapport de transformation \hfill\normalfont\footnotesize E1 sujet 2026 (Q.E10)}
]
\textbf{enonce.}
Pour obtenir $\bar{u}_s = 600\,\mathrm{V}$ avec un pont PD6 ($\bar{u}_s = 0{,}95 \cdot U_M$) alimente par un transformateur etoile-toile (primaire $U_{1,comp} = 400\,\mathrm{V}$), calculer $U_M$, la valeur efficace des tensions composees secondaires $U_{s,comp}$, et le rapport $m$.

\tcbline

\textbf{Corrige.}

\textbf{etape 1 :}
$U_M = 600/0{,}95 \approx 631\,\mathrm{V}$

\textbf{etape 2 :}
$U_{s,comp} = U_M/\sqrt{2} = 631/1{,}414 \approx 446\,\mathrm{V}$

\textbf{etape 3 :}
\[
  m = \frac{U_{s,comp}}{U_{1,comp}} = \frac{446}{400} \approx \mathbf{1{,}12}
\]
\end{tcolorbox}

\vspace{4pt}

% -- QT3 -------------------------------------------------------------

\begin{tcolorbox}[questiontype,
  title={QT3 - PD2 + hacheur - alimentation Urbanloop 230\,V - 72\,V \hfill\normalfont\footnotesize E2 appliquee 2024}
]
\textbf{enonce.}
Le boetier alimentation de la capsule Urbanloop convertit le reseau monophase $230\,\mathrm{V}$ en $72\,\mathrm{V}$ DC sur les rails. La structure retenue est : PD2 + filtre LC + hacheur Buck. Calculer $\bar{U}_{PD2}$ puis le rapport cyclique $\alpha$ du hacheur.

\tcbline

\textbf{Corrige.}

\textbf{etape 1 - tension DC du pont PD2 :}
\[
  \bar{U}_{PD2} = \frac{2\sqrt{2}}{\pi} \times 230 = 0{,}9 \times 230 \approx 207\,\mathrm{V}
\]

\textbf{etape 2 - rapport cyclique du hacheur Buck :}
\[
  \alpha = \frac{\langle u_{sortie} \rangle}{E} = \frac{72}{207} \approx \mathbf{0{,}35}
\]
\textit{Le PD2 fixe la tension DC d'entree ($207\,\mathrm{V}$) ; le hacheur adapte au niveau des rails ($72\,\mathrm{V}$).}
\end{tcolorbox}

\end{tcolorbox}

% -----------------------------------------------------------------------
%  4. ATTENTES DU JURY
% -----------------------------------------------------------------------

\begin{tcolorbox}[
  jury,
  title={4. Attentes du jury - Erreurs recurrentes et points de vigilance},
]

\begin{enumerate}[label=\textbf{\arabic*.}, leftmargin=*, itemsep=5pt]

  \item \textbf{Memoriser les trois formules de $\bar{U}$.}
  PD2 $\approx 0{,}9\,V_{eff}$, P3 $\approx 1{,}17\,V_{s,simple}$, PD6 $\approx 2{,}34\,V_{s,simple}$.
  Les confondre est une erreur eliminatoire dans les questions de dimensionnement.

  \item \textbf{P3 $\neq$ PD6.}
  P3 = 3 diodes (simple alternance, 3 pulsations). PD6 = 6 diodes (pont de Graetz triphase, 6 pulsations). Le terme - PD3 - peut apparaetre dans certains sujets pour designer le PD6.

  \item \textbf{Tension PD6 : utiliser les tensions \emph{composees}.}
  $\bar{U}_{PD6} = 0{,}95\,U_M$ avec $U_M$ = amplitude des tensions \textbf{composees}. Utiliser $V_{simple}$ donne une valeur sous-estimee d'un facteur $\sqrt{3}$.

  \item \textbf{Facteur de puissance $\neq$ $\cos\varphi$.}
  Pour un redresseur - diodes, le courant n'est pas sinusoedal - $f_P \neq \cos\varphi$.
  Utiliser la formule $f_P = \frac{\sqrt{q}}{\pi}\sin\!\left(\frac{\pi}{q}\right)$.

  \item \textbf{Conclure par rapport au cahier des charges.}
  Ne pas juste calculer : comparer $f_P$ - la valeur imposee et conclure (conforme / non conforme).

  \item \textbf{Identifier d'abord la tension DC voulue, puis calculer $m$.}
  Le rapport de transformation du transformateur amont est deduit de la tension secondaire AC necessaire pour obtenir la tension DC cible, pas l'inverse.

\end{enumerate}

\medskip
\begin{tcolorbox}[
  colback=colJuryFrame!12!colJury,
  colframe=colJuryFrame,
  left=4pt, right=4pt, top=2pt, bottom=2pt,
  fontupper=\small\itshape
]
Rapport jury 2025 : "De nombreux candidats ne maetrisent pas les calculs de puissance - distinguer clairement $P$ (active), $Q$ (reactive), $S$ (apparente) dans le bilan du redresseur."\\
Rapport jury 2024 : "Le redresseur PD2 est l'etage d'entree d'un variateur - son rele dans le THD global doit etre identifie ; il injecte des harmoniques de rangs $2k\pm1$ cete reseau."
\end{tcolorbox}

\end{tcolorbox}

% -----------------------------------------------------------------------
%  PIED - PReSENCE DANS LES SUJETS
% -----------------------------------------------------------------------

\begin{tcolorbox}[
  enhanced, breakable,
  colback=colGrayLight,
  colframe=colSeparator,
  left=4pt, right=4pt, top=3pt, bottom=3pt,
  before skip=4pt, after skip=0pt,
  fonttitle=\small\bfseries,
  title={Presence dans les sujets analyses (2013-2026)},
]
\begin{center}
\footnotesize\renewcommand{\arraystretch}{1.15}
\begin{tabular}{lccccccccccccc}
\toprule
\textbf{Annee} & 13 & 14 & 15 & 16 & 17 & 18 & 19 & 20 & 21 & 22 & 23 & 24 & 25 \\
\midrule
E1 (commune)   & $\checkmark$ & $\checkmark$ & $\checkmark$ & - & $\checkmark$ & $\checkmark$ & - & $\checkmark$ & - & - & $\checkmark$ & $\checkmark$ & $\checkmark$ \\
E2 discipl. IE & - & $\checkmark$ & - & $\checkmark$ & - & - & $\checkmark$ & $\checkmark$ & $\checkmark$ & $\checkmark$ & - & $\checkmark$ & $\checkmark$ \\
\bottomrule
\end{tabular}
\end{center}
\smallskip
\footnotesize\textbf{Sous-themes frequents :}
PD6 - tension moyenne et rapport de transfo (7/13) - facteur de puissance (6/13) - PD2 + hacheur (5/13) - harmoniques cete reseau (4/13) - pont commande thyristors (2/13)

\medskip
\textbf{Ordres de grandeur (reperes rapides).}
\begin{itemize}[leftmargin=1.3em]
  \item Reseau BT FR : tension phase-neutre typique \SI{230}{\volt}, tension phase-phase typique \SI{400}{\volt}, frequence \SI{50}{\hertz}.
  \item Courants d'etude frequents en sujet : de quelques \si{\ampere} a quelques dizaines de \si{\ampere}.
  \item En conversion d'energie, un rendement attendu reste generalement entre 0{,}8 et 0{,}98 selon la technologie.
\end{itemize}

\medskip
\textbf{Synthese express (a refaire sans support).}
\begin{enumerate}[leftmargin=1.4em]
  \item Poser les hypotheses du modele et les conventions de signe.
  \item Ecrire 2 a 4 relations cles avec unites.
  \item Realiser un calcul numerique complet et verifier l'ordre de grandeur.
  \item Justifier le resultat par un argument technique de niveau CAPET.
\end{enumerate}

\end{tcolorbox}

\CAPETENDFILE









